\documentclass[10pt, letterpaper]{article}

% Packages:
\usepackage[
    ignoreheadfoot,
    top=2cm,
    bottom=2cm,
    left=2cm,
    right=2cm,
    footskip=1.0cm
]{geometry}
\usepackage{titlesec}
\usepackage{tabularx}
\usepackage{array}
\usepackage[dvipsnames]{xcolor}
\definecolor{primaryColor}{RGB}{0, 0, 0}
\usepackage{enumitem}
\usepackage{fontawesome5}
\usepackage{amsmath}
\usepackage{graphicx}
\usepackage[
    pdftitle={Xinhai Chang's CV},
    pdfauthor={Xinhai Chang},
    pdfcreator={LaTeX with RenderCV},
    colorlinks=true,
    urlcolor=primaryColor
]{hyperref}
\usepackage[pscoord]{eso-pic}
\usepackage{calc}
\usepackage{bookmark}
\usepackage{lastpage}
\usepackage{changepage}
\usepackage{paracol}
\usepackage{ifthen}
\usepackage{needspace}
\usepackage{iftex}

% Ensure ATS parsable:
\ifPDFTeX
    \input{glyphtounicode}
    \pdfgentounicode=1
    \usepackage[T1]{fontenc}
    \usepackage[utf8]{inputenc}
    \usepackage{lmodern}
\fi

% Settings:
\raggedright
\AtBeginEnvironment{adjustwidth}{\partopsep0pt}
\pagestyle{empty}
\setcounter{secnumdepth}{0}
\setlength{\parindent}{0pt}
\setlength{\topskip}{0pt}
\setlength{\columnsep}{0.15cm}
\pagenumbering{gobble}
\titleformat{\section}{\needspace{4\baselineskip}\bfseries\large}{}{0pt}{}[\vspace{1pt}\titlerule]
\titlespacing{\section}{-1pt}{0.3cm}{0.2cm}
\renewcommand\labelitemi{$\vcenter{\hbox{\small$\bullet$}}$}
\newenvironment{highlights}{\begin{itemize}[topsep=0.10cm,parsep=0.10cm,partopsep=0pt,itemsep=0pt,leftmargin=0cm+10pt]}{\end{itemize}}
\newenvironment{highlightsforbulletentries}{\begin{itemize}[topsep=0.10cm,parsep=0.10cm,partopsep=0pt,itemsep=0pt,leftmargin=10pt]}{\end{itemize}}
\newenvironment{onecolentry}{\begin{adjustwidth}{0cm+0.00001cm}{0cm+0.00001cm}}{\end{adjustwidth}}
\newenvironment{twocolentry}[2][]{\onecolentry\def\secondColumn{#2}\setcolumnwidth{\fill, 4.5cm}\begin{paracol}{2}}{\switchcolumn \raggedleft \secondColumn\end{paracol}\endonecolentry}
\newenvironment{header}{\setlength{\topsep}{0pt}\par\kern\topsep\centering\linespread{1.5}}{\par\kern\topsep}

\begin{document}
    \newcommand{\AND}{\unskip\cleaders\copy\ANDbox\hskip\wd\ANDbox\ignorespaces}
    \newsavebox\ANDbox
    \sbox\ANDbox{$|$}

    \begin{header}
        \fontsize{25pt}{25pt}\selectfont XINHAI CHANG
        \vspace{5pt}
        
        \normalsize
        
        Beijing, China \kern 5.0pt \AND \kern 5.0pt
        \href{mailto:changxinhai@stu.pku.edu.cn}{changxinhai@stu.pku.edu.cn} \\
        \vspace{3pt}
        \href{https://chang-xinhai.github.io/}{\faGlobe\ \textbf{chang-xinhai.github.io}} \kern 5.0pt \AND \kern 5.0pt
        \href{https://github.com/chang-xinhai}{\faGithub\ chang-xinhai} \kern 5.0pt \AND \kern 5.0pt
        \href{https://scholar.google.com/citations?user=NVxBzq4AAAAJ}{\faGraduationCap\ Google Scholar}
    \end{header}
    \vspace{5pt - 0.3cm}
    
    \section{Education}

    \begin{twocolentry}{Expected: Jun 2026}
        \textbf{Peking University}, Beijing, China \\
        Bachelor of Science in Data Science, Yuanpei College
    \end{twocolentry}

    \vspace{0.10cm}
    \begin{onecolentry}
        \begin{highlights}
            \item GPA: 3.748/4.0
            \item Coursework: Computer Vision, Machine Learning, Multimodal Learning, Intelligent Robotics [3.9+]
        \end{highlights}
    \end{onecolentry}

    \section{Research Interests}
    \begin{onecolentry}
        My research bridges \textbf{Computer Vision} and \textbf{Robotics} to close the loop of \textbf{Real-Sim-Real} for Embodied AI. I focus on vision-centric \textbf{High-Fidelity 3D Reconstruction and Generation} (Real-to-Sim) to digitize complex physical worlds into simulation-ready assets, and leverage these assets for robotics-centric \textbf{Generalizable Robot Policy Learning} (Sim-to-Real) to deploy robust intelligence in the real world.
    \end{onecolentry}

    \section{Technical Skills}
    \begin{onecolentry}
        \begin{highlights}
            \item \textbf{Robotics and Simulation:} Isaac Lab, Isaac Sim, Curobo, Robot Manipulation, Motion Planning.
            \item \textbf{Languages and Frameworks:} Python, C/C++, Bash, PyTorch, TensorFlow, Pandas, Qt.
            \item \textbf{Algorithms and Vision:} 3D Reconstruction (Gaussian Splatting), Multimodal Learning, 3D Generation.
        \end{highlights}
    \end{onecolentry}

    \section{Research Experience}

    \begin{twocolentry}{Mar 2025 - Present}
        \textbf{CoRe Lab, Peking University} \\
        \textit{Research Assistant (Advised by Prof. Yixin Zhu)}
    \end{twocolentry}
    \vspace{0.10cm}
    \begin{onecolentry}
        \begin{highlights}
            \item Focusing on \textbf{Sim-to-Real} intelligence, specifically large-scale simulation and robust mobile manipulation.
        \end{highlights}
    \end{onecolentry}

    \vspace{0.2cm}

    \begin{twocolentry}{Apr 2024 - Present}
        \textbf{Harvard AI and Robotics Lab} \\
        \textit{Research Assistant (Mentored by Kaichen Zhou)}
    \end{twocolentry}
    \vspace{0.10cm}
    \begin{onecolentry}
        \begin{highlights}
            \item Focusing on \textbf{Real-to-Sim} pipelines, working on high-fidelity 3D reconstruction and digital twin creation.
        \end{highlights}
    \end{onecolentry}

    \section{Publications}
    
    \begin{onecolentry}
        \begin{highlights}
            \item Niu, Y.*, \textbf{Chang, X.}*, Liu, X.*, Jiao, Z., Zhu, Y. (2025). \textit{A Scalable Whole-body Trajectory Generator for Coordinated Mobile Manipulation}. The IEEE/CVF Conference on Computer Vision and Pattern Recognition (CVPR).
            \item Li, H.*, Huang, X.*, \textbf{Chang, X.}*, Zhou, J., Zhao, H. (2025). \textit{TerraX: Visual Terrain Classification Enhanced by Vision-Language Models}. International Conference on Intelligent Robots and Systems (IROS) (Oral).
            \item Zhou, K.*, Wang, Y.*, Chen, G.*, \textbf{Chang, X.}, Beaudouin, G., Zhan, F., Liang, P. P., Wang, M. (2025). \textit{PAGE-4D: Disentangled Pose and Geometry Estimation for 4D Perception}. International Conference on Learning Representations (ICLR).
            \item \textbf{Chang, X.}, Zhou, K. (2026). \textit{Neural Surface Reconstruction from Sparse Views Using Epipolar Geometry}. CVPR 2026 Workshop on Advancements in 2D and 3D Anomaly Detection and Mobile Manipulation Learning (A2AMML).
            \item Tang, C.*, \textbf{Chang, X.}* (2025). \textit{StyleCraft: High-Quality Arbitrary Style Transfer via Unified Content-Style Fusion}. International Conference on Intelligent Computing (ICIC).
        \end{highlights}
    \end{onecolentry}

    \section{Preprints}
    
    \begin{onecolentry}
        \begin{highlights}
            \item Zhou, K.*, \textbf{Chang, X.}*, Kim, T.*, Zhang, J.*, Cao, Y., Peng, C., Zhan, F., Zhao, H., Dong, H., Ting, K. M., Zhu, Y. (2025). \textit{RAD: A Realistic Multi-View Benchmark for Pose-Agnostic Anomaly Detection}. Under Review.
            \item Zhou, K.*, Hong, L.*, \textbf{Chang, X.}*, Zhong, Y., Xie, E., Dong, H., Li, Z., Yang, Y., Li, Z., Zhang, W. (2025). \textit{SplatMesh: Interactive 3D Segmentation and Editing Using Mesh-Based Gaussian Splatting}. Under Review.
        \end{highlights}
        \vspace{0.2cm}
        \footnotesize{* denotes equal contribution}
    \end{onecolentry}
    
    % Ensure these definitions are in your preamble (before \begin{document})
    % or right before the section:
    \definecolor{MacaronRed}{HTML}{FF8B94}
    \definecolor{MacaronBlue}{HTML}{6EB5FF}
    \newcommand{\cvdot}{\textcolor{MacaronRed}{\raisebox{1pt}{\scalebox{0.65}{\faCircle}}}}
    \newcommand{\robdot}{\textcolor{MacaronBlue}{\raisebox{1pt}{\scalebox{0.65}{\faCircle}}}}
    
    % Modified Section Header with Legend
    \section{Academic Projects \hfill \normalfont\normalsize\cvdot\hspace{3pt}\textit{Computer Vision} \hspace{1.5em} \robdot\hspace{3pt}\textit{Robotics}}

    % --- Direction 1: Real-to-Sim ---
    \vspace{0.1cm}
    \begin{onecolentry}
        \textbf{\textit{I. Real-to-Sim: High-Fidelity 3D Reconstruction and Generation}}
    \end{onecolentry}
    \vspace{0.1cm}

    \begin{onecolentry}
        \textbf{EpiS: Neural Surface Reconstruction from Sparse Views Using Epipolar Geometry} \hfill CVPR 2026 Workshop A2AMML
    \end{onecolentry}
    \begin{onecolentry}
        \begin{highlights}
        \item A neural surface reconstruction model resolving geometric ambiguity via epipolar geometry and depth priors.
        \item Responsible for epipolar feature extraction, monocular depth alignment, and baseline evaluation.
        \item Roadmap: Achieves high-fidelity 3D geometry reconstruction from sparse visual observations.\hspace{4pt}\cvdot
        \end{highlights}
    \end{onecolentry}

    \vspace{0.15cm}

    \begin{onecolentry}
        \textbf{SplatMesh: Interactive 3D Segmentation and Editing} \hfill Preprint
    \end{onecolentry}
    \begin{onecolentry}
        \begin{highlights}
        \item An interactive 3D segmentation and editing algorithm integrating Gaussian Splatting with QEM downsampling.
        \item Responsible for QEM downsampling, mesh editing pipeline construction, and model training and evaluation.
        \item Roadmap: Generates editable, memory-efficient digital twins for downstream simulation tasks.\hspace{4pt}\cvdot
        \end{highlights}
    \end{onecolentry}

    \vspace{0.15cm}

    \begin{onecolentry}
        \textbf{PAGE-4D: Disentangled Pose and Geometry Estimation for 4D Perception} \hfill ICLR 2026
    \end{onecolentry}
    \begin{onecolentry}
        \begin{highlights}
        \item A feedforward 4D perception model featuring a dynamics-aware mask to disentangle static and dynamic scenes.
        \item Responsible for novel view synthesis module, LVSM architecture adaptation, and baseline evaluation.
        \item Roadmap: Extends high-fidelity reconstruction to dynamic environments for 4D simulation.\hspace{4pt}\cvdot
        \end{highlights}
    \end{onecolentry}

    \vspace{0.15cm}
    
    \begin{onecolentry}
        \textbf{RAD: Realistic Benchmark for Pose-Agnostic Anomaly Detection} 
        \hfill Preprint
    \end{onecolentry}
    \begin{onecolentry}
        \begin{highlights}
        \item A realistic, robot-captured multi-view benchmark for anomaly detection on reflective and symmetric objects.
        \item Responsible for dataset curation, 2D and 3D model training, and baseline evaluation.
        \item Roadmap: Benchmarks high-fidelity perception against unstructured physical complexities.\hspace{4pt}\cvdot\hspace{2.5pt}\robdot
        \end{highlights}
    \end{onecolentry}

    % --- Direction 2: Sim-to-Real ---
    \vspace{0.3cm}
    \begin{onecolentry}
        \textbf{\textit{II. Sim-to-Real: Generalizable Robot Policy Learning}}
    \end{onecolentry}
    \vspace{0.1cm}

    \begin{onecolentry}
        \textbf{AutoMoMa: Scalable Coordinated Mobile Manipulation Trajectory Generation} 
        \hfill CVPR 2026
    \end{onecolentry}
    \begin{onecolentry}
        \begin{highlights}
        \item A scalable data generation pipeline utilizing GPU-accelerated planning for coordinated mobile manipulation.
        \item Responsible for scene generation, large-scale simulation data collection, and policy training and evaluation.
        \item Roadmap: Generates large-scale manipulation data to enable generalizable policy learning.\hspace{4pt}\robdot
        \end{highlights}
    \end{onecolentry}

    \vspace{0.15cm}

    \begin{onecolentry}
        \textbf{TerraX: Visual Terrain Classification Enhanced by VLMs} 
        \hfill IROS 2025
    \end{onecolentry}
    \begin{onecolentry}
        \begin{highlights}
        \item A vision-language framework combining multimodal data and a CLIP-based VLM for terrain classification.
        \item Responsible for heterogeneous data organization, fine-grained annotation, and baseline evaluation.
        \item Roadmap: Enables robust visual terrain perception for generalizable robotic deployment.\hspace{4pt}\cvdot\hspace{2.5pt}\robdot
        \end{highlights}
    \end{onecolentry}

    
    \section{Honors and Awards}

    \begin{onecolentry}
        \begin{highlights}
            \item \textbf{2025 National Scholarship} (Highest honor for undergraduate students in China)
            \item 2024 Second-class Scholarship, Peking University
            \item First Prize, 21st Jiang Zehan Cup Mathematical Modeling Competition, 2024
            \item Xing Zhengde Scholarship, Yuanpei College, 2023
            \item Third-class Scholarship, Peking University, 2023
        \end{highlights}
    \end{onecolentry}

    \section{Teaching Experience}

    \begin{twocolentry}{Autumn 2025}
        \textbf{Teaching Assistant}, Introduction to Computation B  (Prof. Jun Sun)
    \end{twocolentry}
    \vspace{0.10cm}
    \begin{onecolentry}
        \begin{highlights}
            \item Led weekly programming labs, graded assignments, and designed exam questions.
        \end{highlights}
    \end{onecolentry}

    \vspace{0.2cm}

    \begin{twocolentry}{Autumn 2025}
        \textbf{Teaching Assistant}, Introduction to Computation C (Prof. Baobao Chang)
    \end{twocolentry}
    \vspace{0.10cm}
    \begin{onecolentry}
        \begin{highlights}
            \item Led weekly programming labs, graded assignments, and designed exam questions and the final project.
        \end{highlights}
    \end{onecolentry}

    \vspace{0.2cm}

    \begin{twocolentry}{Autumn 2024}
        \textbf{Teaching Assistant}, Introduction to Computer Systems (Prof. Lu Zhang)
    \end{twocolentry}
    \vspace{0.10cm}
    \begin{onecolentry}
        \begin{highlights}
            \item Led weekly discussion sessions, graded assignments, and contributed to exam design.
        \end{highlights}
    \end{onecolentry}

    \section{Extracurricular Activities}

    \begin{onecolentry}
        \begin{highlights}
            \item Volunteer Teacher, International Asian Liver Center, Gansu, China, Summer 2024
            \item Volunteer, Yunnan Golden Monkey Conservation, Peking University, Winter 2023
        \end{highlights}
    \end{onecolentry}

\end{document}
